

\subsection{Overview}

System integration is a critical part of the Medicare Bot project, where multiple subsystems including the Raspberry Pi, Arduino, backend server, and frontend dashboard are connected to function as a unified system. This integration enables real-time communication, automated decision-making, and synchronized operation across all components.
The integration process builds upon the modules developed in Chapter \ref{softwareimplementation} and the face recognition system discussed in Chapter \ref{facerecognition}.
\subsection{Raspberry Pi and Arduino Integration}

The Raspberry Pi communicates with the Arduino Mega through USB serial communication. The Raspberry Pi acts as the master controller, while the Arduino acts as a slave device responsible for executing physical actions.

\subsubsection{Serial Communication Protocol}

Commands are sent from the Raspberry Pi in the form of encoded strings:

\begin{itemize}
\item \texttt{d1 -> Activate servo 1}
\item \texttt{d2 -> Activate servo 2}
\item \texttt{d3 -> Activate servo 3}
\end{itemize}
Upon receiving these commands, the Arduino interprets them and rotates the corresponding servo motor to dispense medication.

\subsubsection{Challenges and Solutions}

During development, issues related to unstable servo movement and power supply limitations were encountered. These were resolved by:

\begin{itemize}
    \item Separating power supply for Raspberry Pi and servos
    \item Using a dedicated power bank for Arduino
    \item Ensuring stable voltage using a buck converter
\end{itemize}

\subsection{Raspberry Pi and Backend Integration}

The Raspberry Pi communicates with the backend server using HTTP requests. This enables data exchange such as logs, patient information, and images.

\subsubsection{Log Transmission}

After each dispensing event, the Raspberry Pi sends a POST request to the backend API containing:

\begin{itemize}
    \item Patient ID
    \item Dispensing status
    \item Compartment number
    \item Confidence score
    \item Captured image
\end{itemize}

\subsubsection{Image Upload Handling}

Initially, the system encountered a \texttt{422 Unprocessable Entity} error when sending log data. This issue was caused by incorrect request formatting.

The problem was resolved by switching from JSON-based requests to \texttt{multipart/form-data}, allowing both text data and image files to be transmitted correctly.

\subsubsection{API Communication Issues}

Other issues encountered included:

\begin{itemize}
    \item Incorrect API endpoint formatting
    \item Typographical errors in URLs
\end{itemize}

These were resolved through debugging and testing API requests using proper request structures.

\subsection{Backend and Frontend Integration}

The frontend dashboard communicates with the backend server to fetch and display system data.

\subsubsection{Image Display Issue}

A major issue occurred where images were not displayed on the frontend despite being stored correctly.

The root cause was an incorrect URL path:

\begin{verbatim}
/api/static/images -> Incorrect
/static/images -> Correct
\end{verbatim}

Updating the frontend to use the correct path resolved the issue.

\subsubsection{Real-Time Data Updates}

The frontend periodically fetches updated logs from the backend using API calls, enabling near real-time monitoring of system activity.

\subsection{End-to-End System Integration}

The complete system operates through coordinated interaction between all components:

\begin{enumerate}
    \item Raspberry Pi performs face recognition
    \item Medication schedule is verified
    \item Command is sent to Arduino for dispensing
    \item Image is captured
    \item Data is sent to backend server
    \item Backend stores logs and images
    \item Frontend retrieves and displays updated information
\end{enumerate}

\subsection{Integration Challenges}

Several challenges were encountered during system integration:

\begin{itemize}
    \item Power instability affecting servo performance
    \item API communication errors
    \item Incorrect data formats causing request failures
    \item Image path misconfiguration
\end{itemize}

These challenges were resolved through iterative debugging and system testing, leading to a stable and fully integrated system.

\subsection{Current System Status}

After successful integration, the system is capable of:

\begin{itemize}
    \item Real-time face recognition
    \item Automated medication dispensing
    \item Image-based logging
    \item Backend data storage and retrieval
    \item Frontend visualization of logs
\end{itemize}

The system currently operates as a fully functional prototype with all major components successfully integrated.